\begin{frame}{``벨만방정식''이라는 이름의 유래 (역사)}
  이 재귀식 --- ``현재 가치 = 지금 보상 + 할인된 다음 상태
  가치'' --- 을 발견하고 이름을 붙인 사람이 \kb{리처드 벨만
  (Richard Bellman, 1920--1984)}.
  \begin{itemize}
    \item 1950년대 벨만은 ``복잡한 최적화 문제를 작은 부분 문제로
      쪼개서 그 답을 재귀적으로 조합한다''는 \kb{동적계획법}을
      체계화했고, 그 핵심 장치가 바로 이 재귀식.
    \item ``\kb{차원의 저주(curse of dimensionality)}''라는 용어를
      만든 것으로도 유명(이번 장 머리말).
  \end{itemize}
  \vskip0.5em
  \begin{block}{``Bellman equation''은 강화학습만의 개념이 아니다}
    \textbf{최적제어, 금융공학, 확률과정 통계} 등 ``시간에 따라
    상태가 변하며 \textbf{순차적으로 결정}해야 하는 문제'' 전반에서
    쓰이는, 20세기 수학의 \textbf{가장 널리 인용된 재귀식 중 하나}.
    \vskip0.2em
    {\small 강화학습에서 이 방정식이 ``\textbf{기대}''(expectation,
    이 절) 버전과 ``\textbf{최적}''(optimality, $\max$ 이 붙은,
    Week 4) 버전으로 나뉘어 등장하는 것뿐, \textbf{뼈대는 벨만이
    70여 년 전에 세운 것과 정확히 같다}.}
  \end{block}
  \vskip0.3em
  {\footnotesize 이번 절에서는 ``기대'' 버전($\mathbb{E}_\pi$),
  Week 4에서는 ``최적'' 버전($\max_a$)을 배운다.}
\end{frame}
